% ============================================================
%  GUÍA DE PREPARACIÓN - TALLER
%  Cálculo I - Usach Premium
%  Contenidos: límites, continuidad, funciones trigonométricas,
%  funciones exponenciales y logaritmos.
% ============================================================

\documentclass[letterpaper,12pt]{article}

% ============================================================
% PAQUETES
% ============================================================
\usepackage[utf8]{inputenc}
\usepackage[T1]{fontenc}
\usepackage[spanish]{babel}
\usepackage{amsmath,amssymb,amsfonts}
\usepackage{geometry}
\usepackage{fancyhdr}
\usepackage{enumitem}
\usepackage{graphicx}
\usepackage{hyperref}
\usepackage{multicol}
\usepackage{parskip}
\usepackage{xcolor}
\usepackage{tcolorbox}
\usepackage{eso-pic}
\makeatletter
\providecommand{\IfPDFManagementActiveF}[1]{#1}
\makeatother
\usepackage{transparent}
\usepackage{tikz}
\usetikzlibrary{arrows.meta,babel,calc,patterns,shadows.blur}
\usepackage{array}
\usepackage{booktabs}
\usepackage{colortbl}
\usepackage{needspace}
\tcbuselibrary{skins,breakable}

% ============================================================
% GEOMETRÍA Y COLORES
% ============================================================
\geometry{letterpaper,margin=1.5cm,top=3cm,bottom=2.5cm,headheight=2cm}

\definecolor{celesteSuave}{HTML}{F0F7FF}
\definecolor{azulFuerte}{HTML}{1A5276}
\definecolor{azulMedio}{HTML}{2980B9}
\definecolor{turquesaUsach}{HTML}{33CCCC}
\definecolor{enlace}{HTML}{1A5276}
\definecolor{grisTexto}{HTML}{4D5656}
\definecolor{verdeValidado}{HTML}{1E8449}

% ============================================================
% INFORMACIÓN DE LA GUÍA
% ============================================================
\newcommand{\logoup}{../../../../../template/images/logo_up.png}
\newcommand{\logousach}{../../../../../template/images/logo_usach.png}
\newcommand{\logofondo}{../../../../../template/images/logo_up.png}

\newcommand{\institucion}{Usach Premium}
\newcommand{\curso}{Cálculo I}
\newcommand{\guiasnumero}{Guía de preparación - Taller}
\newcommand{\semestre}{1er. Semestre de 2026}
\newcommand{\preparadopor}{Usach Premium.}
\newcommand{\webinstitucion}{www.usachpremium.cl}
\newcommand{\instagraminstitucion}{https://www.instagram.com/usach.premium}
\newcommand{\transparencia}{0.045}
\newcommand{\fraseportada}{"Por y para estudiantes"}

\newcommand{\seno}{\operatorname{sen}}
\newcommand{\Dom}{\operatorname{Dom}}
\newcommand{\Rec}{\operatorname{Rec}}
\newcommand{\R}{\mathbb{R}}

\newcommand{\listacontenidos}{%
    \item[\color{azulFuerte}$\Rightarrow$] Límites y continuidad
    \item[\color{azulFuerte}$\Rightarrow$] Funciones trigonométricas
    \item[\color{azulFuerte}$\Rightarrow$] Funciones exponenciales y logaritmos
    \item[\color{azulFuerte}$\Rightarrow$] 18 ejercicios con solucionario
}

% ============================================================
% HIPERVÍNCULOS, ENCABEZADO Y PIE
% ============================================================
\hypersetup{
    colorlinks=true,
    linkcolor=enlace,
    urlcolor=enlace,
    citecolor=enlace,
    pdfborder={0 0 0},
    pdftitle={Guía de preparación - Taller de Cálculo I},
    pdfauthor={Usach Premium}
}

\pagestyle{fancy}
\fancyhf{}
\fancyhead[L]{\includegraphics[height=1.2cm]{\logoup}}
\fancyhead[R]{%
    \begin{minipage}[b]{0.47\textwidth}
        \raggedleft
        \fontsize{9pt}{11pt}\selectfont
        \textbf{\color{azulFuerte}\institucion}\\
        \curso\ - Preparación Taller\\
        \textit{\color{gray}@usach.premium}
    \end{minipage}\hspace{0.1cm}%
    \includegraphics[height=1.2cm]{\logousach}%
}
\fancyfoot[L]{\fontsize{9pt}{11pt}\selectfont\textit{\fraseportada}}
\fancyfoot[R]{\fontsize{9pt}{11pt}\selectfont Página \thepage}
\renewcommand{\headrulewidth}{0.4pt}
\renewcommand{\footrulewidth}{0.4pt}

\fancypagestyle{plain}{%
    \fancyhf{}
    \renewcommand{\headrulewidth}{0pt}
    \renewcommand{\footrulewidth}{0pt}
}

% ============================================================
% MARCA DE AGUA Y CAJAS
% ============================================================
\newcommand{\MarcaAgua}{%
    \AddToShipoutPicture{%
        \AtPageCenter{%
            \makebox(0,0){%
                \transparent{\transparencia}%
                \includegraphics[width=0.7\textwidth]{\logofondo}%
            }%
        }%
    }%
}

\newtcolorbox{resumenbox}[1]{%
    colback=celesteSuave,
    colframe=azulFuerte,
    colbacktitle=azulFuerte,
    coltitle=white,
    title=\textbf{#1},
    fonttitle=\bfseries,
    arc=7pt,
    boxrule=1pt,
    enhanced,
    breakable
}

\newtcolorbox{ejerciciobox}[1]{%
    colback=white,
    colframe=azulFuerte,
    colbacktitle=azulFuerte,
    coltitle=white,
    title=\textbf{#1},
    fonttitle=\bfseries,
    arc=7pt,
    boxrule=1pt,
    enhanced,
    drop shadow={black!6},
    breakable
}

\newtcolorbox{solbox}[1]{%
    colback=celesteSuave,
    colframe=azulFuerte,
    title=\textbf{#1},
    fonttitle=\bfseries,
    arc=5pt,
    boxrule=0.9pt,
    breakable
}

\newtcolorbox{validacionbox}{%
    colback=verdeValidado!6,
    colframe=verdeValidado,
    coltext=black,
    arc=6pt,
    boxrule=1pt
}

\newcommand{\tituloSeccion}[1]{%
    \needspace{4\baselineskip}
    \subsection*{\color{azulFuerte}#1}
}

\setlist[enumerate]{itemsep=0.4em,topsep=0.4em}

% ============================================================
\begin{document}
\hypersetup{pageanchor=false}
\thispagestyle{plain}

% ============================================================
% PORTADA
% ============================================================
\AddToShipoutPicture*{%
    \put(54,700){\includegraphics[width=2cm]{\logoup}}
    \put(420,706){\includegraphics[width=5cm]{\logousach}}
}

\vspace*{0.5cm}

\begin{center}
    \color{azulFuerte}
    {\huge\textbf{GUÍA DE EJERCICIOS}}\\[0.5cm]
    {\Huge\textbf{\curso}}\\[0.3cm]
    {\Large\guiasnumero}\\[1.5cm]

    \begin{tcolorbox}[
        colback=celesteSuave,
        colframe=azulFuerte,
        arc=10pt,
        width=0.85\textwidth,
        center,
        boxrule=1.5pt]
        \centering
        \vspace{0.2cm}
        {\Large\textbf{Contenidos}}
        \vspace{0.3cm}
        \color{black}
        \begin{itemize}[leftmargin=2.5cm,itemsep=0.35cm]
            \listacontenidos
        \end{itemize}
        \vspace{0.2cm}
    \end{tcolorbox}
\end{center}

\vspace{1.2cm}

\begin{center}
    {\Large\textbf{\semestre}}\\[0.7cm]
    {\large\textbf{\preparadopor}}\\[1cm]
    \begin{tcolorbox}[
        colback=verdeValidado!6,
        colframe=verdeValidado,
        width=0.62\textwidth,
        center,
        arc=6pt,
        boxrule=1pt]
        \centering
        \textcolor{verdeValidado}{\textbf{Solucionario verificado con Python}}
    \end{tcolorbox}
\end{center}

\vfill

\begin{center}
    {\color{azulFuerte}\hrule height 0.5pt}
    \vspace{0.3cm}
    \begin{minipage}{0.45\textwidth}
        \centering
        \textbf{Instagram:}\\
        \small\href{\instagraminstitucion}{@usach.premium}
    \end{minipage}
    \begin{minipage}{0.45\textwidth}
        \centering
        \textbf{Web:}\\
        \small\href{http://\webinstitucion}{\webinstitucion}
    \end{minipage}
\end{center}

\clearpage
\pagenumbering{arabic}
\hypersetup{pageanchor=true}
\MarcaAgua

% ============================================================
% RESUMEN
% ============================================================
\section*{\color{azulFuerte}Antes de comenzar}

\begin{resumenbox}{MAPA RÁPIDO}
\begin{minipage}[t]{0.31\textwidth}
\raggedright
\textbf{\color{azulFuerte}Límites}

Primero sustituye. Si aparece \(0/0\), simplifica mediante factorización o racionalización.

Para que \(f\) sea continua en \(a\), deben coincidir el límite por la izquierda, el límite por la derecha y \(f(a)\).
\end{minipage}\hfill
\begin{minipage}[t]{0.31\textwidth}
\raggedright
\textbf{\color{azulFuerte}Trigonometría}

Si
\[
f(x)=A\seno(Bx)+D,
\]
entonces
\[
\text{amplitud}=|A|,
\]
\[
T=\frac{2\pi}{|B|}.
\]
\[
\seno^2x+\cos^2x=1.
\]
\end{minipage}\hfill
\begin{minipage}[t]{0.31\textwidth}
\raggedright
\textbf{\color{azulFuerte}Exponencial y logaritmo}

\[
b^u=b^v\Longrightarrow u=v
\]
\[
b^x=c\Longleftrightarrow x=\log_b(c)
\]
\[
\log_b(MN)=\log_bM+\log_bN.
\]
Antes de resolver un logaritmo, exige que su argumento sea positivo.
\end{minipage}
\end{resumenbox}

\begin{validacionbox}
\textbf{\color{verdeValidado}Sugerencia de trabajo:}
resuelve primero sin mirar el solucionario. En las ecuaciones, comprueba que cada respuesta pertenezca al dominio y al intervalo pedido.
\end{validacionbox}

% ============================================================
% EJERCICIOS
% ============================================================
\section*{\color{azulFuerte}Ejercicios}

\tituloSeccion{I. Límites y continuidad}
\noindent Calcule los límites de los ejercicios 1 al 5.

\begin{multicols}{2}
\begin{enumerate}[label=\color{azulFuerte}\textbf{\arabic*.},leftmargin=*,itemsep=1.1em]
    \item
    \[
        \lim_{x\to2}\frac{x^2+3x-1}{x+1}.
    \]

    \item
    \[
        \lim_{x\to3}\frac{x^2-9}{x-3}.
    \]

    \item
    \[
        \lim_{x\to4}\frac{\sqrt{x}-2}{x-4}.
    \]

    \item
    \[
        \lim_{x\to0}\frac{\seno(5x)}{x}.
    \]

    \item
    \[
        \lim_{x\to+\infty}
        \frac{3x^2-x+4}{2x^2+5}.
    \]
\end{enumerate}
\end{multicols}

\begin{ejerciciobox}{6. CONTINUIDAD CON PARÁMETROS}
Considere
\[
f(x)=
\begin{cases}
\displaystyle\frac{\sqrt{3x+1}-2}{x-1}, & x<1,\\[2mm]
c, & x=1,\\[2mm]
\displaystyle a\,\frac{\sqrt{x+3}-2}{x-1}, & x>1.
\end{cases}
\]
Determine \(a,c\in\R\) para que \(f\) sea continua en \(x=1\).
\end{ejerciciobox}

\tituloSeccion{II. Funciones trigonométricas}

\begin{multicols}{2}
\begin{enumerate}[label=\color{azulFuerte}\textbf{\arabic*.},leftmargin=*,start=7,itemsep=1em]
    \item Para
    \[
        f(x)=2\seno(3x)-1,
    \]
    determine amplitud, período, eje medio, valor máximo, valor mínimo y recorrido.

    \item Para
    \[
        g(x)=-3\cos\left(\frac{x}{2}\right)+2,
    \]
    determine amplitud, período, eje medio, valor máximo, valor mínimo, recorrido y \(g(0)\).

    \item Resuelva en \(x\in[0,2\pi)\):
    \[
        \seno x=-\frac{\sqrt3}{2}.
    \]

    \item Resuelva en \(x\in[0,2\pi)\):
    \[
        2\cos^2x+3\cos x+1=0.
    \]
\end{enumerate}
\end{multicols}

\begin{ejerciciobox}{11. IDENTIDAD CON ÁNGULO AGUDO}
Sea \(\alpha\) un ángulo agudo que cumple
\[
    \frac{1}{\csc\alpha}=\frac{1+\cos\alpha}{2}.
\]
Determine el valor exacto de \(\seno\alpha\).
\end{ejerciciobox}

\begin{ejerciciobox}{12. MODELO PERIÓDICO}
La altura, en metros, de una boya está dada por
\[
    h(t)=5+4\seno\left(\frac{\pi t}{6}\right),
    \qquad t\geq0,
\]
donde \(t\) se mide en horas. Determine el período, la altura máxima, la altura mínima y el primer instante en que alcanza su altura máxima.
\end{ejerciciobox}

\tituloSeccion{III. Funciones exponenciales y logaritmos}

\begin{multicols}{2}
\begin{enumerate}[label=\color{azulFuerte}\textbf{\arabic*.},leftmargin=*,start=13,itemsep=1em]
    \item Calcule:
    \[
        \log_2(32)+\log_3\left(\frac19\right).
    \]

    \item Resuelva:
    \[
        3^{2x-1}=27.
    \]

    \item Resuelva y revise el dominio:
    \[
        \log_2(x-1)+\log_2(x-3)=3.
    \]

    \item Resuelva:
    \[
        2e^x-5=0.
    \]

    \item Para
    \[
        q(x)=\ln(2x-4)+1,
    \]
    determine dominio, asíntota vertical, recorrido y el punto donde la gráfica corta al eje \(x\).
\end{enumerate}
\end{multicols}

\begin{ejerciciobox}{18. MODELO EXPONENCIAL}
Una población se modela mediante
\[
    P(t)=A\,2^{kt},
\]
donde \(t\) se mide en días. Inicialmente hay \(500\) individuos y después de \(6\) días hay \(2000\).

\begin{enumerate}[label=\alph*),leftmargin=*]
    \item Determine \(A\), \(k\) y la fórmula final de \(P(t)\).
    \item Determine cuándo la población llegará a \(5000\) individuos. Exprese el resultado con logaritmos y entregue también una aproximación decimal.
\end{enumerate}
\end{ejerciciobox}

% ============================================================
% SOLUCIONARIO
% ============================================================
\clearpage
\begin{center}
    \begin{tcolorbox}[
        colback=azulFuerte,
        colframe=azulFuerte,
        colupper=white,
        width=0.8\textwidth,
        center,
        arc=10pt]
        \centering\Large\textbf{SOLUCIONARIO}
    \end{tcolorbox}
\end{center}

\begin{validacionbox}
\centering
\textbf{\color{verdeValidado}18 de 18 resultados verificados con Python y SymPy.}
\end{validacionbox}

\begin{solbox}{I. Límites y continuidad}
\begin{enumerate}[label=\color{azulFuerte}\textbf{\arabic*.},leftmargin=*,itemsep=1em]
    \item Como el denominador no se anula en \(x=2\), sustituimos directamente:
    \[
    \lim_{x\to2}\frac{x^2+3x-1}{x+1}
    =\frac{4+6-1}{3}
    =\boxed{3}.
    \]

    \item Factorizamos \(x^2-9=(x-3)(x+3)\):
    \[
    \lim_{x\to3}\frac{(x-3)(x+3)}{x-3}
    =\lim_{x\to3}(x+3)
    =\boxed{6}.
    \]

    \item Racionalizamos:
    \[
    \frac{\sqrt{x}-2}{x-4}
    \cdot\frac{\sqrt{x}+2}{\sqrt{x}+2}
    =\frac{1}{\sqrt{x}+2}.
    \]
    Por tanto,
    \[
        \lim_{x\to4}\frac{\sqrt{x}-2}{x-4}
        =\boxed{\frac14}.
    \]

    \item Usamos \(\lim_{u\to0}\frac{\seno u}{u}=1\):
    \[
    \lim_{x\to0}\frac{\seno(5x)}{x}
    =5\lim_{x\to0}\frac{\seno(5x)}{5x}
    =\boxed{5}.
    \]

    \item Dividimos numerador y denominador por \(x^2\):
    \[
    \lim_{x\to+\infty}
    \frac{3-\frac1x+\frac4{x^2}}{2+\frac5{x^2}}
    =\boxed{\frac32}.
    \]

    \item Para el tramo izquierdo:
    \[
    \begin{aligned}
    \lim_{x\to1^-}\frac{\sqrt{3x+1}-2}{x-1}
    &=
    \lim_{x\to1^-}
    \frac{3}{\sqrt{3x+1}+2}\\
    &=\frac34.
    \end{aligned}
    \]
    Para el tramo derecho:
    \[
    \begin{aligned}
    \lim_{x\to1^+}
    a\frac{\sqrt{x+3}-2}{x-1}
    &=
    a\lim_{x\to1^+}
    \frac{1}{\sqrt{x+3}+2}\\
    &=\frac{a}{4}.
    \end{aligned}
    \]
    La continuidad exige
    \[
        \frac34=c=\frac{a}{4}.
    \]
    Así,
    \[
        \boxed{a=3,\qquad c=\frac34}.
    \]
\end{enumerate}
\end{solbox}

\begin{solbox}{II. Funciones trigonométricas}
\begin{enumerate}[label=\color{azulFuerte}\textbf{\arabic*.},leftmargin=*,start=7,itemsep=1em]
    \item En \(f(x)=2\seno(3x)-1\):
    \[
    \text{amplitud}=2,
    \qquad
    \text{período}=\frac{2\pi}{3},
    \qquad
    \text{eje medio}:y=-1.
    \]
    Como \(-1\leq\seno(3x)\leq1\),
    \[
        \boxed{\max f=1,\quad \min f=-3,\quad \Rec(f)=[-3,1]}.
    \]

    \item En \(g(x)=-3\cos(x/2)+2\):
    \[
    \text{amplitud}=3,
    \qquad
    \text{período}=4\pi,
    \qquad
    \text{eje medio}:y=2.
    \]
    Por tanto,
    \[
        \boxed{\max g=5,\quad \min g=-1,\quad \Rec(g)=[-1,5]}.
    \]
    Además,
    \[
        \boxed{g(0)=-3\cos0+2=-1}.
    \]

    \item El seno es negativo en los cuadrantes III y IV, con ángulo de referencia \(\pi/3\):
    \[
        \boxed{x\in\left\{\frac{4\pi}{3},\frac{5\pi}{3}\right\}}.
    \]

    \item Factorizamos como una ecuación cuadrática en \(\cos x\):
    \[
    2\cos^2x+3\cos x+1
    =(2\cos x+1)(\cos x+1)=0.
    \]
    Entonces \(\cos x=-\frac12\) o \(\cos x=-1\), de donde
    \[
        \boxed{x\in\left\{\frac{2\pi}{3},\pi,\frac{4\pi}{3}\right\}}.
    \]

    \item Como \(1/\csc\alpha=\seno\alpha\),
    \[
        2\seno\alpha=1+\cos\alpha
        \quad\Longrightarrow\quad
        \cos\alpha=2\seno\alpha-1.
    \]
    Elevamos al cuadrado y usamos \(\cos^2\alpha=1-\seno^2\alpha\):
    \[
    1-\seno^2\alpha=(2\seno\alpha-1)^2
    \Longrightarrow
    \seno\alpha(5\seno\alpha-4)=0.
    \]
    Como \(\alpha\) es agudo,
    \[
        \boxed{\seno\alpha=\frac45}.
    \]

    \item El argumento del seno aumenta en \(2\pi\) cuando \(t\) aumenta en \(12\), por lo que el período es
    \[
        \boxed{T=12\text{ horas}}.
    \]
    Como \(-1\leq\seno(\pi t/6)\leq1\),
    \[
        \boxed{h_{\min}=1\text{ m},\qquad h_{\max}=9\text{ m}}.
    \]
    El primer máximo ocurre cuando
    \[
        \frac{\pi t}{6}=\frac{\pi}{2},
    \]
    así que
    \[
        \boxed{t=3\text{ horas}}.
    \]
\end{enumerate}
\end{solbox}

\begin{solbox}{III. Funciones exponenciales y logaritmos}
\begin{enumerate}[label=\color{azulFuerte}\textbf{\arabic*.},leftmargin=*,start=13,itemsep=1em]
    \item
    \[
    \log_2(32)+\log_3\left(\frac19\right)
    =5+(-2)
    =\boxed{3}.
    \]

    \item Como \(27=3^3\),
    \[
    3^{2x-1}=3^3
    \Longrightarrow
    2x-1=3
    \Longrightarrow
    \boxed{x=2}.
    \]

    \item El dominio exige
    \[
        x-1>0,\qquad x-3>0
        \quad\Longrightarrow\quad x>3.
    \]
    Usamos la propiedad del producto:
    \[
    \log_2\bigl((x-1)(x-3)\bigr)=3
    \Longrightarrow
    (x-1)(x-3)=8.
    \]
    Así,
    \[
    x^2-4x-5=0
    \Longrightarrow
    (x-5)(x+1)=0.
    \]
    Se descarta \(x=-1\) por no pertenecer al dominio. Entonces,
    \[
        \boxed{x=5}.
    \]

    \item
    \[
    2e^x-5=0
    \Longrightarrow
    e^x=\frac52
    \Longrightarrow
    \boxed{x=\ln\left(\frac52\right)}.
    \]

    \item El argumento del logaritmo debe ser positivo:
    \[
        2x-4>0
        \Longrightarrow
        \boxed{\Dom(q)=(2,+\infty)}.
    \]
    Cuando \(x\to2^+\), \(\ln(2x-4)\to-\infty\), por lo que la asíntota vertical es
    \[
        \boxed{x=2}.
    \]
    El recorrido de un logaritmo trasladado sigue siendo
    \[
        \boxed{\Rec(q)=\R}.
    \]
    Para hallar el corte con el eje \(x\):
    \[
    \ln(2x-4)+1=0
    \Longrightarrow
    2x-4=e^{-1}.
    \]
    Por tanto, el punto es
    \[
        \boxed{\left(2+\frac{1}{2e},\,0\right)}.
    \]

    \item De \(P(0)=500\):
    \[
        A\,2^{k\cdot0}=A=500.
    \]
    Luego, usando \(P(6)=2000\):
    \[
    500\,2^{6k}=2000
    \Longrightarrow
    2^{6k}=4=2^2
    \Longrightarrow
    k=\frac13.
    \]
    El modelo es
    \[
        \boxed{P(t)=500\,2^{t/3}}.
    \]
    Para llegar a \(5000\):
    \[
    500\,2^{t/3}=5000
    \Longrightarrow
    2^{t/3}=10
    \Longrightarrow
    \frac{t}{3}=\log_2(10).
    \]
    Finalmente,
    \[
        \boxed{t=3\log_2(10)\approx9{,}97\text{ días}}.
    \]
\end{enumerate}
\end{solbox}

\vspace{0.8cm}
\begin{center}
    \small\textbf{Usach Premium} - Nos vemos en la ayudantía.\\
    \textit{"Por y para estudiantes."}
\end{center}

\end{document}
